\chapter{Stage 6: The Calibrated Composite and Four-Outcome Gate}
\label{ch:composite}

The verifier of \Cref{ch:verifier} produced nine numbers. Stage 6 turns those nine
numbers into one, the stored trust score $\ustored$, and then turns that one number
into an action. This is the chapter where the whole architecture's promise comes
due: the number it produces must be an \emph{honest probability}, so that a single
fixed threshold can mean the same thing on a fact-check, a trivia question, and a
multiple-choice item. We build the number from the ground up, teaching every piece,
and then we open the gate that reads it.

\begin{intuition}[Two jobs, in order]
Stage 6 has two jobs. The first is \emph{fusion}: take nine signals that live on
different scales and point in different directions and combine them into one
calibrated probability that the answer is correct. The second is \emph{decision}:
compare that probability against fixed cut-points and commit to one of four actions,
store, defer, abstain, or discard. The fusion is where the cleverness lives; the
decision is deliberately simple, because all the hard calibration work has already
been done by the time the gate reads the number.
\end{intuition}

\section{Why fusion is hard: scales and sign-flips}
\label{sec:composite-problem}

The naive way to combine nine signals is a weighted average. It fails here for a
concrete reason that is worth seeing before we fix it.

\begin{intuition}[The same signal can mean opposite things on different benchmarks]
Consider the question-answer relevance signal $\qarel$, which measures how well the
answer addresses the question. On an open-ended trivia question, a high relevance is
good news: an answer that squarely addresses the question is more likely correct. On
a fact-checking claim, where the ``answer'' is a verdict like \emph{supports} or
\emph{refutes}, a very high surface relevance can actually be \emph{bad} news,
because the model is echoing the claim's framing rather than judging it. Measured on
the calibration data, this signal correlates \emph{negatively} with correctness on
the fact-checking benchmark and \emph{positively} on the trivia benchmark. A single
fixed weight, positive or negative, must be wrong on one of them. No weighted
average can serve both.
\end{intuition}

\begin{bythenumbers}[The sign-flip, measured]
On the cold-start calibration fold, the relevance signal $\qarel$ correlates with
correctness at a Pearson coefficient of about $-0.275$ on FEVER, a moderate
\emph{negative} relation, and about $+0.068$ on TriviaQA, a \emph{positive} one. The
sign itself differs between benchmarks, so the isotonic fit of
\Cref{sec:composite-isotonic} learns a \emph{decreasing} curve on FEVER and an
\emph{increasing} one on TriviaQA. A single fixed weight, positive or negative, must
be wrong on one of them. The flip is not even stable over time: by the third cycle,
after the generator had been fine-tuned, the measured FEVER relation had drifted to
weakly positive, which is one reason the calibration is re-fitted every cycle
(\Cref{sec:composite-training}). This single measurement rules out any shared
fixed-weight combination and motivates the per-signal, per-benchmark calibration that
follows.
\end{bythenumbers}

The composite solves this in three steps, drawn in \Cref{fig:composite-pipeline}.
First, each signal is passed through its own \emph{calibration curve} that converts
its raw value into a probability of correctness, learning the right direction
automatically. Second, those curves are fitted separately per benchmark and gently
blended toward a shared curve to avoid overfitting. Third, the per-signal
probabilities are combined into one score through a learned weighting. We take each
step in turn.

\begin{figure}[t]
\centering
\resizebox{\textwidth}{!}{%
\begin{tikzpicture}[
  font=\small\sffamily,
  sig/.style={rectangle, rounded corners=3pt, draw=cbScenario!70!black, fill=cbScenario!10,
    line width=0.9pt, align=center, inner sep=4pt, text width=2.0cm, font=\footnotesize,
    minimum height=0.8cm},
  stp/.style={rectangle, rounded corners=4pt, line width=1pt, align=center,
    inner sep=5pt, text width=2.7cm, minimum height=1.3cm,
    drop shadow={shadow xshift=0.3ex, shadow yshift=-0.3ex, fill=black!10}},
  num/.style={rectangle, rounded corners=4pt, draw=cbNumbers!75!black, fill=cbNumbers!12,
    line width=1.1pt, align=center, inner sep=5pt, text width=1.8cm, minimum height=1.0cm,
    drop shadow={shadow xshift=0.3ex, shadow yshift=-0.3ex, fill=black!10}},
  flow/.style={-{Stealth[length=2.4mm]}, line width=1pt, draw=cbInk!75},
]
  \node[sig] (s) at (0,0) {\textbf{nine signals}\\{\footnotesize raw values}};
  \node[stp, draw=cbIntuition!75!black, fill=cbIntuition!10] (iso) at (3.4,0)
    {\textbf{1. per-signal}\\\textbf{isotonic}\\{\footnotesize raw $\to P_i$}};
  \node[stp, draw=cbFormula!70!black, fill=cbFormula!10] (shr) at (7.0,0)
    {\textbf{2. per-benchmark}\\\textbf{+ shrinkage}\\{\footnotesize $\alpha=0.6$ toward pooled}};
  \node[stp, draw=cbExample!70!black, fill=cbExample!10] (boost) at (10.6,0)
    {\textbf{3. log-odds sum}\\\textbf{+ boost}\\{\footnotesize $\sigma(\sum w_i \ell_i + b)$}};
  \node[num] (u) at (13.9,0) {$\ustored$\\{\footnotesize $\in[0,1]$}};
  \draw[flow] (s) -- (iso);
  \draw[flow] (iso) -- (shr);
  \draw[flow] (shr) -- (boost);
  \draw[flow] (boost) -- (u);
\end{tikzpicture}}
\caption[The composite pipeline]{The three-step composite. Each raw signal is first
mapped to a calibrated probability $P_i$ by its own monotone curve; the curves are
fitted per benchmark and shrunk toward a pooled curve; and the per-signal
probabilities are combined in log-odds space through a learned weighting and a
sigmoid into the single trust score $\ustored$. (Schematic.)}
\label{fig:composite-pipeline}
\end{figure}

\section{Step one: turning one raw signal into a calibrated probability}
\label{sec:composite-isotonic}

The first step solves a small version of the whole problem: take one signal, say the
top-passage entailment, and convert its raw value into a probability that the answer
is correct. The tool is \emph{isotonic regression}, and because the rest of the
chapter rests on it, we teach it from scratch.

\begin{intuition}[What ``calibrated'' means here]
A raw signal is just a number on some scale; its value of $0.8$ has no inherent
meaning. We want to replace it with a probability: ``answers for which this signal
reads $0.8$ turn out correct seventy percent of the time, so map $0.8 \mapsto 0.70$.''
A function that does this faithfully, so that the output really is the empirical
fraction correct, is a \emph{calibration curve}. Once every signal has been put
through its own calibration curve, all nine speak the same language: probability of
correctness.
\end{intuition}

\subsection*{Isotonic regression and the monotone step function}

To learn the curve we use a labelled \emph{calibration fold}: a held-out set of
answers for which we know the raw signal value and the ground-truth correctness (one
if the answer was right, zero if wrong). Plotting correctness against the signal
gives a cloud of zeros and ones. We want a curve through that cloud, and we impose
one assumption: the curve must be \emph{monotone}, that is, always non-decreasing (or
always non-increasing). The assumption is principled, since a well-designed signal
should not reverse direction in the middle of its range, and it is exactly what makes
the fit data-efficient.

\begin{precise}[Isotonic regression, the pool-adjacent-violators idea]
Isotonic regression finds the monotone function that best fits the labelled points
in the least-squares sense. The classic algorithm is intuitive: sort the points by
signal value, walk left to right, and wherever the running fitted values would
violate monotonicity (a later average dips below an earlier one), \emph{pool} those
adjacent points and replace them with their shared average. Repeated, this produces a
\emph{step function}: a piecewise-constant curve that jumps up at a set of breakpoints
called \emph{knots} and is flat between them. \caem{} stores each signal's curve as
its list of knot positions and knot heights, and at run time reads off a value by
linear interpolation between the two surrounding knots. The output is clipped to the
range $[0.02,\,0.98]$ so that no single signal can claim certainty, and a signal needs
at least fifty calibration points to earn a curve at all; below that it is dropped and
contributes nothing.
\end{precise}

\begin{figure}[t]
\centering
\begin{tikzpicture}
\begin{axis}[
  width=10.5cm, height=6.2cm,
  xlabel={raw $\pgmax$ (top-passage grounding)},
  ylabel={$P(\text{correct})$},
  xmin=0, xmax=1, ymin=-0.08, ymax=1.08,
  xtick={0,0.25,0.5,0.75,1}, ytick={0,0.25,0.5,0.75,1},
  axis lines=left, font=\small\sffamily, tick label style={font=\scriptsize},
  legend style={at={(0.03,0.97)}, anchor=north west, draw=cbInk!25, font=\scriptsize, fill=white},
  grid=both, grid style={cbInk!8},
]
  % schematic 0/1 cloud (the underlying fold points are not stored; shown for intuition)
  \addplot[only marks, mark=*, mark size=1.0pt, cbPitfall!45] coordinates {
    (0.02,0.02)(0.05,0.0)(0.09,0.03)(0.15,0.0)(0.22,0.02)(0.31,0.0)(0.40,0.03)(0.55,0.0)(0.66,0.02)(0.04,0.0)(0.12,0.02)(0.27,0.0)};
  \addlegendentry{incorrect ($\mathrm{em}=0$), schematic}
  \addplot[only marks, mark=*, mark size=1.0pt, cbExample!50] coordinates {
    (0.39,1.0)(0.50,0.98)(0.66,1.0)(0.78,0.98)(0.87,1.0)(0.93,0.98)(0.96,1.0)(0.97,0.98)(0.99,1.0)(0.70,1.0)(0.85,0.98)(0.60,1.0)};
  \addlegendentry{correct ($\mathrm{em}=1$), schematic}
  % REAL deployed isotonic step function for p_ground_max (global, cycle 3)
  \addplot[const plot, cbIntuition, line width=1.6pt] coordinates {
    (0,0.385)(0.006,0.486)(0.049,0.500)(0.069,0.510)(0.153,0.539)(0.254,0.575)(0.369,0.583)(0.389,0.638)(0.664,0.670)(0.867,0.721)(0.961,0.778)(0.969,0.923)(0.988,0.929)};
  \addlegendentry{deployed $\pgmax$ calibration}
  % read-off example
  \draw[dashed, cbNumbers!80!black] (axis cs:0.50,0) -- (axis cs:0.50,0.638) -- (axis cs:0,0.638);
  \node[font=\scriptsize, cbNumbers!70!black, anchor=west] at (axis cs:0.04,0.70) {$0.50 \mapsto 0.64$};
\end{axis}
\end{tikzpicture}
\caption[An isotonic calibration curve]{The blue step function is the \emph{actual
deployed} calibration curve for the top-passage grounding signal $\pgmax$, read
directly from the cycle-three calibration file. Isotonic regression fitted this
monotone curve from the labelled calibration fold and \caem{} reads it by linear
interpolation: a raw grounding of $0.50$ maps to a calibrated
$P(\text{correct}) \approx 0.64$, and a raw value near $1.0$ maps to about $0.93$.
The curve never reaches the corners because the output is clamped into $[0.02, 0.98]$,
so no single signal asserts certainty. The scattered points are schematic, drawn only
to show the labelled cloud the curve is fitted through; the per-sample fold values are
not stored in the calibration file.}
\label{fig:isotonic}
\end{figure}

\subsection*{Learning the direction automatically}

The sign-flip of \Cref{sec:composite-problem} is handled here, for free. Before
fitting, the procedure measures the correlation between the raw signal and
correctness; if it is positive it fits a non-decreasing curve, and if it is negative
it fits a non-increasing one. No hand-coded table of per-signal directions is needed.
\Cref{fig:signflip} shows the same relevance signal calibrated in opposite directions
on the two benchmarks.

\begin{figure}[t]
\centering
\begin{tikzpicture}
\begin{axis}[
  name=l, width=6.9cm, height=5.2cm,
  title={\footnotesize FEVER: relevance reads \emph{down} ($r=-0.275$)},
  xlabel={$\qarel$ (raw)}, ylabel={$P(\text{correct})$},
  xmin=0.5,xmax=0.75,ymin=0,ymax=1, xtick={0.5,0.6,0.7}, ytick={0,0.5,1},
  axis lines=left, font=\small\sffamily, tick label style={font=\scriptsize}, grid=both, grid style={cbInk!8},
]
  % REAL cold-start FEVER q_a_relevance knots (decreasing)
  \addplot[cbPitfall, line width=1.5pt, const plot] coordinates {
    (0.567,0.793)(0.582,0.627)(0.665,0.634)(0.666,0.604)(0.678,0.591)(0.682,0.578)(0.687,0.558)(0.689,0.435)(0.705,0.412)(0.726,0.514)};
\end{axis}
\begin{axis}[
  name=r, xshift=7.7cm, width=6.9cm, height=5.2cm,
  title={\footnotesize TriviaQA: relevance reads \emph{up} ($r=+0.068$)},
  xlabel={$\qarel$ (raw)},
  xmin=0.5,xmax=0.75,ymin=0,ymax=1, xtick={0.5,0.6,0.7}, ytick={0,0.5,1},
  axis lines=left, font=\small\sffamily, tick label style={font=\scriptsize}, grid=both, grid style={cbInk!8},
]
  % REAL cold-start TriviaQA q_a_relevance knots (increasing)
  \addplot[cbExample!70!black, line width=1.5pt, const plot] coordinates {
    (0.500,0.020)(0.5005,0.085)(0.5052,0.214)(0.7082,0.335)(0.7098,0.585)(0.7311,0.720)};
\end{axis}
\end{tikzpicture}
\caption[The same signal, calibrated both ways]{The two curves are the \emph{actual
deployed} relevance calibrations from the cold-start calibration file, fitted
independently per benchmark. On FEVER (left) higher relevance correlates with
\emph{lower} correctness (Pearson $r=-0.275$), so isotonic learns a decreasing curve;
on TriviaQA (right) it correlates with \emph{higher} correctness ($r=+0.068$), so the
curve increases. The relevance score itself lives in a narrow band near $0.5$ to
$0.73$, which is why the axes are zoomed. The auto-detected direction is what lets one
signal serve both benchmarks honestly; the jaggedness of the FEVER curve, fitted on
only a few hundred points, is exactly what the shrinkage of the next section tames.}
\label{fig:signflip}
\end{figure}

\begin{pitfall}[Why isotonic, and not the simpler temperature or histogram]
Three reasons, each grounded in the report's design rationale. A single-parameter
fit like the temperature scaling of \Cref{ch:upre} can only soften or sharpen; it
cannot bend, and several signals (notably entailment, which saturates near the top of
its range) need a genuine bend. A histogram, which just bins the signal and reports
each bin's hit rate, wastes the small calibration fold by giving every bin its own
independent count; isotonic shares information across neighbouring values through its
monotone structure and so needs far less data. And isotonic reads its direction off
the data, which is precisely what the sign-flip demands. Temperature scaling is right
for the one-dimensional pre-routing confidence; isotonic is right for these nine.
\end{pitfall}

\section{Step two: per-benchmark curves, shrunk toward a shared one}
\label{sec:composite-shrinkage}

Fitting one curve per signal across all data would average away the sign-flip. So the
composite fits a curve per signal \emph{per benchmark}. But a per-benchmark fold is
small, only a few hundred answers, and a flexible curve fitted on little data
\emph{overfits}: it chases noise and generalises worse than a coarser fit. The
composite controls this with \emph{shrinkage}.

\begin{precise}[James-Stein shrinkage toward the pooled curve]
Two curves are fitted for each signal: a \emph{pooled} curve on all benchmarks
together, and a \emph{per-benchmark} curve on each benchmark's slice alone. The
deployed curve is a convex blend of the two, taken at every knot,
\[
  \text{curve}^{b}(x) \;=\; \alpha \cdot \text{curve}^{b}_{\text{local}}(x)
  \;+\; (1-\alpha)\cdot \text{curve}^{\text{pooled}}(x),
  \qquad \alpha = 0.6.
\]
At $\alpha=1$ the deployed curve is purely local (maximally adapted, maximally
overfit); at $\alpha=0$ it is the pooled curve (no per-benchmark adaptation). The
locked value $\alpha=0.6$ keeps a moderate per-benchmark adaptation while pulling each
local curve most of the way back toward the shared one. This is the James-Stein idea:
an estimate from little data is improved by shrinking it toward a pooled prior. The
pooled curve also serves as the fallback for any query whose benchmark is unknown at
deployment.
\end{precise}

\begin{bythenumbers}[Why $\alpha = 0.6$ was chosen]
The locked value comes from a cycle-zero diagnostic. Fitting pure per-benchmark
curves ($\alpha = 1$) made the weak-signal benchmarks \emph{worse}: their ability to
separate correct from incorrect answers, measured by the area under the
receiver-operating curve of the stored trust score, fell by roughly five to six
points below the pooled fit (from about 0.645 to 0.592 on the misconception probe and
from about 0.573 to 0.512 on the multi-step-reasoning probe), because roughly five
hundred answers per benchmark is too little for the flexible local curve. Blending at $\alpha = 0.6$ moved those readings back to parity with the pooled
fit while preserving a real gain on the strong-signal benchmarks. Shrinkage buys
robustness on the benchmarks that need it without surrendering adaptation on the ones
that can afford it.
\end{bythenumbers}

\begin{figure}[t]
\centering
\begin{tikzpicture}
\begin{axis}[
  width=10.5cm, height=5.8cm,
  xlabel={raw $\pgmax$ (grounding) on the commonsense fold}, ylabel={$P(\text{correct})$},
  xmin=0,xmax=1, ymin=0.3,ymax=1, xtick={0,0.25,0.5,0.75,1}, ytick={0.5,0.75,1},
  axis lines=left, font=\small\sffamily, tick label style={font=\scriptsize},
  legend style={at={(0.03,0.97)}, anchor=north west, draw=cbInk!25, font=\scriptsize, fill=white},
  grid=both, grid style={cbInk!8},
]
  % REAL reconstructed pure-local commonsense curve (alpha=1): local=(deployed-0.4*pooled)/0.6
  \addplot[cbPitfall!60, line width=1.3pt, const plot, dashed] coordinates {
    (0.006,0.528)(0.016,0.664)(0.062,0.699)(0.219,0.733)(0.945,0.733)(0.949,0.800)(0.969,0.980)(0.984,0.980)};
  \addlegendentry{per-benchmark, $\alpha{=}1$ (overfit)}
  % REAL pooled (global) p_ground_max curve
  \addplot[cbInk!55, line width=1.3pt, const plot] coordinates {
    (0,0.385)(0.05,0.500)(0.15,0.539)(0.37,0.583)(0.39,0.638)(0.66,0.670)(0.87,0.721)(0.97,0.923)(0.99,0.929)};
  \addlegendentry{pooled (smooth prior)}
  % REAL deployed commonsense blend (alpha=0.6)
  \addplot[cbExample!65!black, line width=1.8pt, const plot] coordinates {
    (0.006,0.511)(0.016,0.595)(0.062,0.619)(0.219,0.655)(0.945,0.728)(0.949,0.769)(0.969,0.957)(0.984,0.959)};
  \addlegendentry{blended, $\alpha{=}0.6$ (deployed)}
\end{axis}
\end{tikzpicture}
\caption[Shrinkage toward the pooled curve]{All three curves are real, for the
grounding signal $\pgmax$ on the commonsense benchmark. The pure per-benchmark fit
(red, dashed, reconstructed from the deployed blend) climbs steeply: on that small
fold grounding looks far more informative than it robustly is, since its true
correlation there is only about $+0.095$. The pooled curve (grey), fitted across all
benchmarks, rises gently and is well supported. The deployed curve (green) is their
$\alpha=0.6$ blend, sitting between the two and damping the local curve's
over-eagerness back toward the pooled prior.}
\label{fig:shrinkage}
\end{figure}

\section{Step three: combining nine probabilities into one}
\label{sec:composite-boost}

After steps one and two, each of the nine signals has produced a calibrated
probability $P_i$ that the answer is correct. We must fuse the nine into one. Doing
this well needs one more idea, the \emph{log-odds}, which we teach now.

\begin{intuition}[Log-odds: the natural currency for combining evidence]
A probability $P$ lives in $[0,1]$, which is awkward for adding evidence: two pieces
of strong evidence cannot push past $1$. The \emph{odds} $P/(1-P)$ live in
$[0,\infty)$, and the \emph{log-odds} $\ell = \log\!\big(P/(1-P)\big)$ live on the
whole real line, with $\ell = 0$ meaning ``even chance'' ($P = 0.5$), positive
meaning ``likely correct,'' and negative ``likely wrong.'' Log-odds are the right
currency because independent pieces of evidence \emph{add} in log-odds space, so
accumulating evidence becomes plain addition rather than an awkward combination of
bounded fractions. So the composite converts
each $P_i$ to its log-odds $\ell_i$, adds them up, and converts the total back to a
probability with the sigmoid $\sigma(z) = 1/(1+e^{-z})$, the inverse of the log-odds
map.
\end{intuition}

\begin{precise}[The log-odds sum, with a learned boost]
The composite score is
\[
  \ustored \;=\; \sigma\!\Big(\, b \;+\; \sum_{i=1}^{9} w_i\,\ell_i \,\Big),
  \qquad
  \ell_i = \log\!\frac{P_i}{1-P_i},
\]
where $P_i$ is the calibrated probability from signal $i$. With every weight $w_i = 1$
and intercept $b = 0$ this is the plain log-odds sum, which treats the nine signals as
independent evidence. \caem{} goes one step further and \emph{learns} the weights
$w_i$ and intercept $b$, which is the \emph{boost layer}.
\end{precise}

\subsection*{The boost layer is a logistic regression}

The boost layer is a logistic regression, and since the term has not yet been taught,
here it is in full.

\begin{intuition}[What a logistic regression is, and how it is trained]
A logistic regression predicts a probability from a set of input features. It forms a
weighted sum of the features plus a constant, $z = b + \sum_i w_i x_i$, and squashes
that sum through the sigmoid to get a probability $\sigma(z)$. \emph{Training} it
means choosing the weights $w_i$ and the constant $b$ that make those predicted
probabilities match the known labels as closely as possible, measured by the
\emph{log-loss} (the negative log-likelihood, which rewards confident-correct
predictions and punishes confident-wrong ones). The fit is found by standard convex
optimisation; there is one global best answer. Here the features $x_i$ are the nine
per-signal log-odds $\ell_i$, and the labels are the calibration fold's correctness
flags, so the boost learns \emph{how much to trust each signal} relative to the
others.
\end{intuition}

\begin{precise}[Regularisation, and why it is set very strong]
Left unconstrained, a logistic regression on correlated features can learn large,
unstable weights that overfit the calibration fold. \emph{Regularisation} penalises
large weights; its strength is set by a constant $C$, where a \emph{smaller} $C$ means
a \emph{stronger} penalty and weights kept closer to zero. \caem{} fits the boost with
$C = 0.01$, a deliberately strong penalty. The reason is that the per-signal isotonic
curves have already done the heavy calibration; the boost is only a gentle correction
for residual correlations among signals (for instance, the two internal-confidence
signals partly measure the same thing). Strong regularisation keeps the boost close to
the plain log-odds sum, correcting only what the data clearly support and refusing to
chase noise on a small fold.
\end{precise}

\begin{figure}[t]
\centering
\begin{tikzpicture}
\begin{axis}[
  width=12.0cm, height=7.0cm,
  xbar, y axis line style={opacity=0}, axis x line=bottom, x axis line style={cbInk!40},
  xlabel={contribution to the log-odds sum (toward correct $\rightarrow$)},
  xmin=-0.32, xmax=0.40, enlarge y limits=0.07, bar width=8pt,
  ytick=data, font=\small\sffamily, tick label style={font=\scriptsize},
  symbolic y coords={base rate, internal, ground-atomic, ground-mean, token, ground-max, dropout, s-avg, q-relevance, entailment},
  nodes near coords, nodes near coords style={font=\tiny, /pgf/number format/fixed, /pgf/number format/precision=3},
  xmajorgrids, grid style={cbInk!8},
]
  \addplot[draw=cbExample!70!black, fill=cbExample!22] coordinates {
    (0.310,entailment) (0.024,q-relevance) (0.010,s-avg) (0.004,dropout) (0.003,ground-max)
    (-0.001,token) (-0.014,ground-mean) (-0.014,ground-atomic) (-0.054,internal) (-0.237,base rate)
  };
  \draw[cbInk!50, dashed] (axis cs:0,base rate) -- (axis cs:0,entailment);
\end{axis}
\end{tikzpicture}
\caption[Signals voting in log-odds space]{The \emph{actual} per-signal log-odds
decomposition for a real FEVER claim, ``Drake Bell is a model,'' computed from the
deployed cycle-three calibration. The top-passage entailment is the one strong vote
that the answer is correct ($+0.310$), but it is almost entirely cancelled by the
FEVER base-rate prior carried in the boost intercept ($-0.237$) and the low
internal-confidence signal ($-0.054$); grounding lands near neutral because the
retrieved passage barely supports the claim. The ten contributions sum to a
hair-positive $+0.032$, and the sigmoid maps that to $\ustored = 0.508$, which the
gate reads as a deferral rather than a store.}
\label{fig:logodds}
\end{figure}

\section{How the composite is fitted: the training phase}
\label{sec:composite-training}

Everything so far has described how the composite \emph{runs}. But nothing in it is
hand-set: every knot height of every isotonic curve, and every weight of the boost
layer, is \emph{learned from data}. That learning happens in a phase kept strictly
separate from the phase in which the composite is used, and the separation is the
key to understanding when and how the composite is trained.

\begin{intuition}[Two phases: fit-time and predict-time]
The composite has a \emph{fit-time} and a \emph{predict-time}, and they never
overlap. At fit-time the curves and weights are \emph{estimated} from a batch of
labelled examples, the way any supervised model is trained. At predict-time, when a
real query arrives, those curves and weights are \emph{frozen}: a signal value is
simply looked up on its stored curve and the stored weights are simply applied. No
fitting happens while answering a query. The composite is trained in batches and
applied one query at a time, exactly mirroring the two clocks of the whole system in
\Cref{ch:bigpicture}.
\end{intuition}

\subsection*{The training data: a labelled calibration fold}

The composite is fitted on a \emph{calibration fold}: a held-out set of roughly
fifteen hundred answers, each one carrying two things, the nine raw signal values the
verifier computed for it, and a \emph{correctness label}, written $\mathrm{em}$, that
is one if the answer was actually right and zero if it was wrong (on TruthfulQA the
correctness is decided by the language-model judge of \Cref{ch:benchmarks} rather than
by string match). In machine-learning terms the nine signals are the \emph{inputs} and
$\mathrm{em}$ is the \emph{target}: the composite learns the mapping from signals to
the probability that $\mathrm{em} = 1$. The fold is held out in the strict sense that
the generator never trains on it; it exists only to teach the calibration what each
signal's values are worth.

\subsection*{Fitting the curves, then the weights}

Fitting proceeds in the same three steps the composite later runs, but now estimating
rather than applying.

\begin{enumerate}[leftmargin=1.6em, itemsep=2pt]
  \item \textbf{Each isotonic curve is fitted} by taking that signal's
  (value,\,$\mathrm{em}$) pairs across the fold and running the
  pool-adjacent-violators procedure of \Cref{sec:composite-isotonic} to get the
  monotone step function and its knots. The direction is read from the sign of the
  signal-to-correctness correlation, and a signal with fewer than fifty usable pairs
  is dropped. This is done per benchmark and per pooled union, then blended by the
  shrinkage of \Cref{sec:composite-shrinkage}.
  \item \textbf{The boost weights are fitted} by building a table whose rows are the
  fold's answers and whose columns are that answer's nine per-signal calibrated
  log-odds (computed by passing each signal through the curve just fitted). A logistic
  regression is then trained on this table against the $\mathrm{em}$ labels, choosing
  the weights and intercept that minimise the log-loss under the strong regularisation
  $C = 0.01$. At least fifty complete rows are required, or the boost is skipped and
  the plain log-odds sum is used.
  \item \textbf{The result is serialised}: the curves (as knot lists) and the boost
  weights are written to a calibration file, which is all that predict-time needs.
\end{enumerate}

\begin{lstlisting}[caption={Fitting the composite, condensed from \texttt{caem/verification/cal\_prob\_composite.py}.}, label={lst:fit}]
for sig in signals:                                  # one isotonic curve per signal
    vals, labels = collect(fold, sig)                 # (signal value, em) pairs from the fold
    if len(vals) < 50: continue                       # too little data -> drop the signal
    direction = "increasing" if corr(vals, labels) >= 0 else "decreasing"  # auto-direction
    curve[sig] = isotonic_fit(vals, labels, direction)   # pool-adjacent-violators -> knots

X = [[log_odds(curve[s].predict(row[s])) for s in signals] for row in fold]  # per-signal log-odds
y = [row["em"] for row in fold]                       # correctness labels
boost = LogisticRegression(C=0.01).fit(X, y)          # minimise log-loss -> weights + intercept
save(curve, boost.coef_, boost.intercept_)            # serialise to composite_calibration.json
\end{lstlisting}

\subsection*{When it is fitted: once at cold-start, then every cycle}

The composite is fitted for the first time at \emph{cold-start} (cycle zero), on the
cycle-zero calibration fold, and the file it produces is the one the system ships with.
It is then \emph{re-fitted at every cycle boundary that commits}, on the same fold
\emph{re-scored under the updated model}. \Cref{fig:composite-train} draws the two
lanes.

\begin{figure}[t]
\centering
\resizebox{\textwidth}{!}{%
\begin{tikzpicture}[
  font=\small\sffamily,
  lane/.style={rectangle, rounded corners=6pt, line width=1pt, inner sep=4pt},
  tb/.style={rectangle, rounded corners=3pt, draw=cbFormula!70!black, fill=cbFormula!10,
    line width=0.9pt, align=center, inner sep=4pt, text width=2.6cm, font=\footnotesize, minimum height=0.9cm},
  ab/.style={rectangle, rounded corners=3pt, draw=cbExample!70!black, fill=cbExample!10,
    line width=0.9pt, align=center, inner sep=4pt, text width=2.6cm, font=\footnotesize, minimum height=0.9cm},
  store/.style={rectangle, rounded corners=3pt, draw=cbNumbers!75!black, fill=cbNumbers!12,
    line width=1pt, align=center, inner sep=4pt, text width=2.2cm, font=\footnotesize, minimum height=0.9cm},
  flow/.style={-{Stealth[length=2.4mm]}, line width=1pt, draw=cbInk!75},
]
  % TRAIN lane
  \node[font=\footnotesize\bfseries, text=cbFormula!75!black, anchor=west] at (-0.4,2.6)
    {FIT-TIME \;(cold-start, then each committed cycle)};
  \node[tb] (fold) at (0.6,1.7) {labelled fold\\{\footnotesize 9 signals + em, $n{\approx}1500$}};
  \node[tb] (fiso) at (4.0,1.7) {fit isotonic\\{\footnotesize per signal}};
  \node[tb] (fboost) at (7.2,1.7) {fit boost\\{\footnotesize logistic reg., $C{=}0.01$}};
  \node[store] (json) at (10.4,1.7) {calibration\\file\\{\footnotesize knots + weights}};
  \draw[flow] (fold) -- (fiso);
  \draw[flow] (fiso) -- (fboost);
  \draw[flow] (fboost) -- (json);
  % APPLY lane
  \node[font=\footnotesize\bfseries, text=cbExample!55!black, anchor=west] at (-0.4,-0.4)
    {PREDICT-TIME \;(every query, frozen)};
  \node[ab] (sig) at (0.6,-1.3) {query's\\9 signals};
  \node[ab] (apply) at (4.0,-1.3) {interpolate curves\\+ weighted log-odds};
  \node[ab] (sigm) at (7.2,-1.3) {sigmoid};
  \node[store, draw=cbNumbers!75!black, fill=cbNumbers!14] (u) at (10.4,-1.3) {$\ustored$};
  \draw[flow] (sig) -- (apply);
  \draw[flow] (apply) -- (sigm);
  \draw[flow] (sigm) -- (u);
  % JSON feeds apply
  \draw[flow, draw=cbInk!55, dashed] (json.south) -- ++(0,-0.5) -| (apply.north)
    node[pos=0.25, right, font=\scriptsize, text=cbInk!60] {load frozen curves + weights};
\end{tikzpicture}}
\caption[Fit-time versus predict-time]{The two phases. At fit-time (top) the labelled
calibration fold is used to estimate the isotonic curves and then the boost weights,
which are saved to a calibration file; this happens once at cold-start and again at
every committed cycle boundary. At predict-time (bottom) a real query's nine signals
are run through the \emph{frozen} curves and weights to produce $\ustored$. The
training never touches a live query, and a live query never re-trains the composite. (Schematic.)}
\label{fig:composite-train}
\end{figure}

\begin{precise}[Why the fold is fixed but re-scored each cycle]
The calibration fold is the \emph{same} fifteen hundred questions every cycle; it is
not resampled. What changes is the \emph{model that answers them}. Because the
self-improvement loop of \Cref{ch:cycle} fine-tunes the generator each cycle, the same
question now produces different signal values and possibly a different correctness
label, so the fold is \emph{re-scored} under the updated model before the composite is
re-fitted on it. Re-fitting is necessary precisely because the model drifts: if the
curves stayed frozen while the signal distributions shifted underneath them, a stored
score of $0.60$ would slowly stop meaning a sixty-percent chance of correctness, and
the fixed gate thresholds would lose their meaning. Holding the fold fixed while
re-scoring it ensures that any change in the fitted calibration reflects the model's
movement, not noise from a different sample. The composite and the pre-routing
temperature are, in fact, the only parts of the architecture re-fitted each cycle;
everything else in the configuration is locked.
\end{precise}

\begin{pitfall}[The composite is supervised; the signals it calibrates are not its labels]
A common confusion is to think the verifier's signals \emph{are} the training labels.
They are not. The signals are the \emph{inputs}; the single training label per answer
is its ground-truth correctness $\mathrm{em}$, obtained by scoring the fold's answers
against gold. The composite is an ordinary supervised model learning ``given these
nine signal readings, how often was the answer actually right?'' That is why a
\emph{labelled} fold is required, and why the quality of $\ustored$ ultimately rests
on having honest correctness labels on the calibration fold.
\end{pitfall}

\section{The whole computation, end to end}
\label{sec:composite-worked}

\begin{workedexample}[A borderline FEVER claim becomes $\ustored = 0.508$]
Take a real FEVER claim the deployed system actually scored, ``Drake Bell is a
model.'' Its nine raw signals enter step one and each is mapped through its own
FEVER-calibrated, shrunk curve to a probability of correctness: the top-passage
entailment reads moderately high (the retrieved passage does engage the claim), the
relevance maps just above one half, the grounding signals land almost exactly at one
half because the passage barely supports the specific assertion, and the
internal-confidence signal maps below one half because the model was unsure. Step
three converts each to log-odds and forms the boosted weighted sum, drawn in
\Cref{fig:logodds}: the single strong upward vote from entailment is almost entirely
cancelled by the FEVER base-rate prior carried in the intercept and by the downward
internal-confidence vote, leaving a hair-positive total of $+0.032$. The sigmoid turns
that into $\ustored = 0.508$. The number is honest: it says this answer is just above
an even chance of being correct, which is exactly the regime where the system should
hesitate rather than commit. The gate reads it as a deferral, and the claim is parked
for reconsideration at the next cycle boundary.
\end{workedexample}

\begin{lstlisting}[caption={The composite score, condensed from \texttt{caem/verification/cal\_prob\_composite.py}.}, label={lst:composite}]
log_odds = boost_intercept                       # b (0 when boost not fitted)
for sig, curve in calibrations.items():          # the nine per-signal curves
    p = curve.predict(signal_values[sig])         # raw value -> calibrated P_i (isotonic)
    p = clip(p, 0.02, 0.98)                        # bound any single signal
    ell = log(p / (1 - p))                         # P_i -> log-odds
    log_odds += boost_weights.get(sig, 1.0) * ell  # weighted (boost) sum
u_stored = 1.0 / (1.0 + exp(-log_odds))            # sigmoid back to [0, 1]
\end{lstlisting}

\begin{fromcode}[Per-benchmark dispatch, and the cold-start fallback]
At run time the composite is handed the query's benchmark tag. If a per-benchmark
calibration exists for that benchmark, the score is computed through its curves and
weights; otherwise the pooled fallback is used. On the very first cycle, before any
calibration fold has been scored, no curves exist yet, so the verifier bootstraps with
a simple fixed weighted sum and switches to the calibrated composite as soon as the
cycle-zero fold is fitted. Thereafter the curves are refit at every cycle boundary
(\Cref{ch:cycle}), which is what keeps $\ustored$ meaning the same thing as the model
underneath it changes.
\end{fromcode}

\section{The four-outcome gate}
\label{sec:composite-gate}

Now the second job. The gate reads $\ustored$ and commits to one of four actions. It
is deliberately a short cascade of fixed thresholds, because the number it reads is
already a calibrated probability.

\begin{figure}[t]
\centering
\begin{tikzpicture}[
  font=\small\sffamily,
  dec/.style={diamond, aspect=2.2, draw=cbFormula!75!black, fill=cbFormula!10,
    line width=1pt, align=center, inner sep=1pt, font=\footnotesize, text width=2.5cm},
  act/.style={rectangle, rounded corners=4pt, line width=1pt, align=center,
    inner sep=4pt, font=\footnotesize, minimum height=0.8cm, text width=2.2cm,
    drop shadow={shadow xshift=0.3ex, shadow yshift=-0.3ex, fill=black!10}},
  flow/.style={-{Stealth[length=2.2mm]}, line width=1pt, draw=cbInk!75},
]
  \node[dec, draw=cbPitfall!75!black, fill=cbPitfall!10] (ee) at (0,3.2)
    {early-exit?\\{\scriptsize $\uint{\ge}0.70 \wedge \pgmax{\le}0.20$}};
  \node[dec] (st) at (0,1.4) {$\ustored \ge 0.60$?};
  \node[dec] (df) at (0,-0.4) {$\ustored \ge 0.45$?};
  \node[dec] (ab) at (0,-2.2) {$\pgmax < 0.20$?};

  \node[act, draw=cbPitfall!70!black, fill=cbPitfall!12] (disc1) at (4.3,3.2) {\textbf{DISCARD}\\{\scriptsize + regenerate}};
  \node[act, draw=cbExample!70!black, fill=cbExample!14] (store) at (4.3,1.4) {\textbf{STORE}\\{\scriptsize into memory}};
  \node[act, draw=cbNumbers!75!black, fill=cbNumbers!14] (defer) at (4.3,-0.4) {\textbf{DEFER}\\{\scriptsize buffer, recheck}};
  \node[act, draw=cbIntuition!75!black, fill=cbIntuition!12] (abst) at (4.3,-2.2) {\textbf{ABSTAIN}\\{\scriptsize refuse honestly}};
  \node[act, draw=cbInk!55, fill=cbInk!6] (disc2) at (4.3,-3.7) {\textbf{DISCARD}\\{\scriptsize drop silently}};

  \draw[flow] (ee) -- node[above,font=\scriptsize,text=cbPitfall!70!black]{yes} (disc1);
  \draw[flow] (ee) -- node[right,font=\scriptsize]{no} (st);
  \draw[flow] (st) -- node[above,font=\scriptsize,text=cbExample!55!black]{yes} (store);
  \draw[flow] (st) -- node[right,font=\scriptsize]{no} (df);
  \draw[flow] (df) -- node[above,font=\scriptsize,text=cbNumbers!65!black]{yes} (defer);
  \draw[flow] (df) -- node[right,font=\scriptsize]{no} (ab);
  \draw[flow] (ab) -- node[above,font=\scriptsize,text=cbIntuition!65!black]{yes} (abst);
  \draw[flow] (ab) -- node[right,font=\scriptsize]{no} (disc2);
\end{tikzpicture}
\caption[The four-outcome gate]{The decision cascade. The confabulation early-exit of
\Cref{ch:verifier} is checked first and forces a discard-and-regenerate when the model
is internally sure but ungrounded. Otherwise the calibrated score is compared against
the storage cut-point $0.60$ and the deferral cut-point $0.45$; below deferral, an
answer with essentially no grounding is turned into an honest abstention, and anything
else is discarded silently. (Schematic.)}
\label{fig:gate}
\end{figure}

\begin{precise}[The decision tree]
After the early-exit, the gate is four lines, read top to bottom, first match wins:
\begin{itemize}[leftmargin=1.4em, itemsep=1pt]
  \item $\ustored \ge 0.60 \Rightarrow$ \textbf{store} the answer as a verified episode.
  \item $0.45 \le \ustored < 0.60 \Rightarrow$ \textbf{defer}: hold it in a buffer for
  reconsideration at the next cycle boundary.
  \item $\ustored < 0.45$ and $\pgmax < 0.20 \Rightarrow$ \textbf{abstain}: the answer
  is both low-confidence and ungrounded, so return an honest refusal.
  \item otherwise $\Rightarrow$ \textbf{discard}: drop the answer silently.
\end{itemize}
The storage and deferral cut-points are \emph{precision targets}, read as ``store only
when the calibrated chance of correctness is at least sixty percent.'' Because
$\ustored$ is a calibrated probability, that sentence means the same thing on every
benchmark.
\end{precise}

\begin{lstlisting}[caption={The gate, condensed from \texttt{caem/verification/verifier.py}.}, label={lst:gate}]
if u_stored >= 0.60:          # store_threshold
    return "STORE"
if u_stored >= 0.45:          # defer_threshold
    return "DEFERRED"
if p_ground_max < 0.20:       # abstain_pground_ceiling
    return "ABSTAIN"          # low confidence AND ungrounded -> honest refusal
return "DISCARD"
\end{lstlisting}

\begin{intuition}[Why four outcomes and not two]
A plain accept/reject gate would throw away two distinctions the architecture depends
on. The first is \emph{defer}: an answer just below the storage line is promising but
not yet trusted, so rather than discard it the system parks it and lets a later, better
model re-judge it. This is exactly how the FEVER claim at $\ustored = 0.508$ is later
promoted to memory (\Cref{ch:cycle}). The second is the split between \emph{abstain}
and \emph{discard} among the low-confidence answers. If the answer is also ungrounded
(its best supporting passage barely entails it, $\pgmax < 0.20$), the honest move is to
tell the user ``I do not know'' rather than emit a guess; that is abstain. If instead
there is some grounding but the composite still rejects it, the answer is dropped
silently as a discard. The four outcomes are the difference between a system that
merely filters and one that knows the difference between ``wrong'' and ``unknown.''
\end{intuition}

\section{Why one fixed threshold governs five benchmarks}
\label{sec:composite-fixed}

The deepest payoff of the chapter is also the simplest to state. The gate uses one
pair of cut-points, $0.60$ and $0.45$, for every benchmark. That is only sound because
all the benchmark-specific work happened upstream: each signal was calibrated per
benchmark, and the composite produced a probability whose meaning is benchmark
independent. A storage score of $0.60$ certifies the same sixty-percent calibrated
chance of correctness whether the query was a fact-check or a trivia question. Push the
calibration into the signals and the threshold can be a constant; skip the calibration
and the threshold would have to be retuned per benchmark, which is exactly the
brittleness the composite exists to remove.

\begin{defence}[If an examiner asks: ``wasn't there a conformal-prediction gate here?'']
An earlier design used a split-conformal gate, which offers a formal coverage
guarantee by calibrating a cut-point on a held-out fold. It was removed. The
guarantee relies on the calibration fold and the evaluation stream being
\emph{exchangeable}, a technical condition meaning their order carries no information,
and a cycle-zero diagnostic found that condition violated in this setting: the
self-improvement loop shifts the answer distribution between the fold and the stream,
so the coverage guarantee no longer held. Rather than ship a guarantee that does not
hold, the design reverted to the fixed-threshold gate above, whose meaning is
maintained instead by re-fitting the calibration every cycle. Reporting the removal is
the honest choice; a guarantee that fails its premise is worse than an explicit
threshold that does not claim one.
\end{defence}

\begin{bythenumbers}[Composite and gate constants]
\begin{itemize}[leftmargin=1.3em, itemsep=1pt]
  \item Composite signals: nine (the verifier's other computed signals do not enter).
  \item Per-signal calibration: isotonic, output clipped to $[0.02, 0.98]$, minimum $50$ samples.
  \item Shrinkage toward pooled curve: $\alpha = 0.6$.
  \item Boost: logistic regression on per-signal log-odds, regularisation $C = 0.01$.
  \item Combination: $\ustored = \sigma\big(b + \sum_i w_i \ell_i\big)$.
  \item Gate: store $\ge 0.60$, defer $\ge 0.45$, abstain if below $0.45$ and $\pgmax < 0.20$, else discard.
  \item Early-exit (first): $\uint \ge 0.70$ and $\pgmax \le 0.20$ forces discard-and-regenerate.
\end{itemize}
\end{bythenumbers}

The composite has now turned nine signals into one honest probability, and the gate
has turned that probability into one of four actions. \Cref{ch:output} takes up
Stage 7, the output contract: how each of the four decisions is rendered to the user,
including how an abstention becomes a calibrated refusal and how the served answer
carries its confidence with it.
